% Laget automatisk av elevark/lag_elevark.py fra sim/stromhuske.html. Endringer her blir overskrevet.
\documentclass[11pt]{article}
\usepackage{elevark}
\begin{document}
\elevarkhode{Fysikk 2 · Magnetiske felt og krefter}{Strømhuske}{https://lektorbo.no/sim/stromhuske.html}{Simuleringen}

En metallstav henger i to lette snorer i et loddrett magnetfelt. Når det går strøm i staven, får den en vannrett kraft $F = IlB$, og huska svinger ut til siden. Her ser du staven fra enden, så strømmen går inn i eller ut av planet.

\section*{Oppgaver}

Start med $I = 2{,}0\,\mathrm{A}$ inn i planet, $B = 0{,}40\,\mathrm{T}$ opp, $l = 10\,\mathrm{cm}$ og $m = 40\,\mathrm{g}$.

\begin{oppgaveliste}
\oppgave \fase{gjett}{Gjett}Bruk høyrehåndsregelen. Svinger huska mot høyre eller venstre? Hva skjer hvis du snur strømmen? Og hvis du snur både strømmen og feltet?
\svar
\oppgave \fase{test}{Regn}Regn ut den magnetiske kraften og tyngden. Vis at $\tan  α = IlB / mg$, og regn ut vinkelen. Sjekk med simuleringen.
\svar
\oppgave \fase{gjett}{Gjett}Hva skjer med vinkelen hvis du dobler strømmen? Hva hvis du i stedet dobler både lengden og massen til staven? Test begge.
\svar
\oppgave \fase{test}{Regn}Hvor stor må strømmen være for at snorene skal danne 30° med loddlinjen? Test svaret.
\svar
\oppgave \fase{test}{Mål}Velg «Ukjent (regn det ut)». Les av vinkelen, og regn ut magnetfeltet. Sjekk med «Vis feltet».
\svar
\oppgave \fase{forklar}{Forklar}Trykk på «Slipp fra loddrett». Huska svinger forbi likevektsvinkelen før den legger seg i ro. Omtrent hvor stor blir den største vinkelen, sammenlignet med vinkelen i likevekt? Hvorfor?
\svar
\oppgave \fase{utfordring}{Utfordring}Tenk deg at magnetfeltet pekte langs staven i stedet for loddrett. Hva ville skjedd med huska? Hva om feltet var vannrett, men vinkelrett på staven?
\svar
\end{oppgaveliste}

\end{document}
